\begin{prob}\textbf{Actual Midterm 2, October 30, Question 3.} On $\mathbb R^3$, let $T$ have standard matrix $A=\left(\begin{smallmatrix}2&1&0\\0&2&0\\0&0&-1\end{smallmatrix}\right)$. (a) Determine its eigenvalues and corresponding eigenspaces. (b) Determine its minimal polynomial with justification. (c) Is $T$ diagonalizable? Explain.\\\emph{Hint:} Use the shifted-kernel equations and test both the quadratic and cubic products; roots alone do not justify the exponent.
\end{prob}
\begin{prob}\textbf{Practice Midterm 2, Question 3.} Let $T$ act on a nonzero finite-dimensional complex space. You may use eigenvalue existence. Prove upper-triangularization by induction: (a) handle dimension $1$; (b) choose an eigenvalue $\lambda$, put $U=\im(T-\lambda I)$, and prove invariance and $\dim U<\dim V$; (c) obtain a basis $u_1,\ldots,u_m$ of $U$ with $Tu_j\in\operatorname{span}(u_1,\ldots,u_j)$; (d) extend it by $v_1,\ldots,v_r$ and use $Tv_j=(T-\lambda I)v_j+\lambda v_j$.\\\emph{Hint:} Proposition 61 explains why the range is the useful smaller space; include the case $U=\{0\}$.
\end{prob}
\begin{prob}\textbf{Practice Midterm 2, Question 4.} Let $T$ act on a nonzero finite-dimensional real space. (a) If $m_T(0)\ne0$, prove invertibility by expressing $T^{-1}$ as a polynomial. (b) If $m_T(0)=0$, prove noninvertibility using minimality. (c) For $R(x,y)=(-y,x)$ on $\mathbb R^2$, find $m_R$. (d) Use (a) to express $R^{-1}$ in terms of $R$ and check directly.\\\emph{Hint:} Factor out $T$ from every nonconstant term; in (b), an inverse would cancel a factor $x$ and produce a smaller relation.
\end{prob}
\begin{prob}\textbf{Midterm 2 Review, Problem 4.} (1) On the real space $V=\operatorname{span}(1,\sin x,\cos x,\sin2x,\cos2x)$, find the eigenvalues of differentiation and decide diagonalizability. (2) For $T(x,y)=(y,-x)$ on $\mathbb C^2$, find eigenvalues and a basis giving an upper-triangular matrix. (3) Show the same rule on $\mathbb R^2$ has no upper-triangular basis.\\\emph{Hint:} Write the derivative equations in the displayed function basis; in (2),(3), check the field and the sign of the rotation.
\end{prob}
